\documentclass[11pt,respuestas,a4]{aleph-examen}
% \documentclass[10pt,a5]{aleph-examen}

% -- Paquetes extra 
\usepackage{aleph-comandos}
\usepackage{systeme}

% -- Datos
\institucion{Facultad de Ciencias Exactas, Naturales y Ambientales}
\carrera{Catálogo STEM}
\asignatura{Matemática III}
\tema{Examen Escrito 1}
\autor{Andrés Merino}
\fecha{Periodo 2026-1}

\logouno[0.16\textwidth]{Logos/logoPUCE_04_ac}
\definecolor{colortext}{HTML}{1957d1}
\definecolor{colordef}{HTML}{1957d1}
\fuente{montserrat}


\begin{document}

\encabezado

\vspace{-5mm}
\section*{Indicaciones}
\begin{itemize}[leftmargin=*]
\item
    En esta actividad se evalúa si el estudiante (Criterio 2.2) calcula derivadas de funciones de varias variables.
\item
    Los ejercicios deben ser resueltos a manoy entregarse en hojas de forma física al docente, sin ayuda de resúmenes o asistentes computacionales.
\item
    Se encuentra prohibido el uso de cualquier otra fuente de información durante todo el examen.
\item
    No está permitido tener ningún dispositivo electrónico encendido durante el examen, incluyendo celulares, tablets, computadoras, audífonos, relojes inteligentes, ni similares.
\item
    En caso de considerar que existe un error en la pregunta o que esta se encuentra mal planteada, se debe indicar cuál es el error y justificarlo.
\item
    Todas las soluciones deben estar correctamente redactadas y explicadas.
\end{itemize}

\section*{Ejercicios}

\begin{preguntas}

%%%%%%%%%%%%%%%%%%%%%%%%%%%%%%%%%%%%%%%%%%
\item
Considere la función
\[
    f(x,y) = \begin{cases}
        \dfrac{(x-1)^2(y-2)}{(x-1)^4 + (y-2)^2} & \text{si } (x,y)\neq(1,2),\\[6pt]
        0 & \text{si } (x,y)=(1,2).
    \end{cases}
\]
\begin{enumerate}[leftmargin=*]
    \item Calcule los límites iterados cuando $(x,y)\to(1,2)$.\puntaje{3}
    \item Calcule el límite por la trayectoria $y=2x$.\puntaje{2}
    \item Calcule el límite cuando $(x,y)\to(1,2)$ por la trayectoria $y-2=(x-1)^2$.\puntaje{2}
    \item ¿Es $f$ continua en $(1,2)$? Justifique a partir de los incisos anteriores.\puntaje{3}
\end{enumerate}


\begin{respuesta}
\begin{enumerate}
    \item Calculemos los límites iterados. Primero cuando $y\to 2$ y luego cuando $x\to 1$:
    \[
        \lim_{x\to 1}\left(\lim_{y\to 2}\frac{(x-1)^2(y-2)}{(x-1)^4+(y-2)^2}\right)
        =
        \lim_{x\to 1}\frac{(x-1)^2(0)}{(x-1)^4+0}
        =
        \lim_{x\to 1} 0
        =
        0.
    \]
    Ahora cuando $x\to 1$ y luego cuando $y\to 2$:
    \[
        \lim_{y\to 2}\left(\lim_{x\to 1}\frac{(x-1)^2(y-2)}{(x-1)^4+(y-2)^2}\right)
        =
        \lim_{y\to 2}\frac{(0)^2(y-2)}{(0)^4+(y-2)^2}
        =
        \lim_{y\to 2} 0
        =
        0.
    \]

    \item Consideremos la trayectoria $y=2x$, es decir, $y-2=2(x-1)$. Entonces:
    \[
        f(x,2x)
        =
        \frac{(x-1)^2\cdot 2(x-1)}{(x-1)^4+4(x-1)^2}
        =
        \frac{2(x-1)^3}{(x-1)^2\bigl((x-1)^2+4\bigr)}
        =
        \frac{2(x-1)}{(x-1)^2+4}.
    \]
    Por tanto:
    \[
        \lim_{x\to 1} f(x,2x)
        =
        \lim_{x\to 1}\frac{2(x-1)}{(x-1)^2+4}
        =
        \frac{0}{4}
        =
        0.
    \]

    \item Consideremos la trayectoria $y-2=(x-1)^2$. Entonces:
    \[
        f\bigl(x,\,2+(x-1)^2\bigr)
        =
        \frac{(x-1)^2\cdot(x-1)^2}{(x-1)^4+\bigl((x-1)^2\bigr)^2}
        =
        \frac{(x-1)^4}{(x-1)^4+(x-1)^4}
        =
        \frac{(x-1)^4}{2(x-1)^4}
        =
        \frac{1}{2}.
    \]
    Por tanto:
    \[
        \lim_{x\to 1} f\bigl(x,\,2+(x-1)^2\bigr)
        =
        \frac{1}{2}.
    \]

    \item Como los límites por distintas trayectorias son diferentes (por $y=2x$ se obtiene $0$, pero por $y-2=(x-1)^2$ se obtiene $\tfrac{1}{2}$), se concluye que:
    \[
        \lim_{(x,y)\to(1,2)} f(x,y)
    \]
    no existe. Por lo tanto, $f$ no es continua en $(1,2)$.\qedhere
\end{enumerate}
\end{respuesta}


%%%%%%%%%%%%%%%%%%%%%%%%%%%%%%%%%%%%%%%%%%
\item
Sea $f(x,y)=x\cos(y^2)\cdot e^{x^2}$.
\begin{enumerate}[leftmargin=*]
    \item Calcule $\dfrac{\partial f}{\partial x}$ y $\dfrac{\partial f}{\partial y}$.\puntaje{4}
    \item Calcule $\nabla f(P)$ en el punto $P=(1,0)$.\puntaje{2}
    \item Calcule la derivada direccional de $f$ en $P$ en la dirección del vector $\overrightarrow{PQ}$, donde $Q=(3,1)$.\puntaje{2}
    \item Calcule la derivada direccional máxima de $f$ en $P$ e indique la dirección en la que se alcanza.\puntaje{2}
\end{enumerate}


\begin{respuesta}
\begin{enumerate}
    \item Calculemos las derivadas parciales usando la regla del producto y la regla de la cadena:
    \[
        \frac{\partial f}{\partial x}(x,y)
        =
        \cos(y^2)\cdot e^{x^2}+x\cos(y^2)\cdot 2x e^{x^2}
        =
        \cos(y^2)\cdot e^{x^2}(1+2x^2).
    \]
    \[
        \frac{\partial f}{\partial y}(x,y)
        =
        x\cdot(-\sen(y^2))\cdot 2y\cdot e^{x^2}
        =
        -2xy\,\sen(y^2)\cdot e^{x^2}.
    \]

    \item Evaluamos en $P=(1,0)$:
    \[
        \frac{\partial f}{\partial x}(1,0)
        =
        \cos(0)\cdot e^{1}\cdot(1+2)
        =
        3e,
        \qquad
        \frac{\partial f}{\partial y}(1,0)
        =
        -2(1)(0)\,\sen(0)\cdot e^{1}
        =
        0.
    \]
    Por lo tanto:
    \[
        \nabla f(P)
        =
        (3e,\ 0).
    \]

    \item Calculamos el vector $\overrightarrow{PQ}$:
    \[
        u = \overrightarrow{PQ}
        =
        Q-P
        =
        (3,1)-(1,0)
        =
        (2,1).
    \]
    La derivada direccional de $f$ en $P$ en la dirección de $u$ es:
    \[
        f'(P;u)
        =
        \nabla f(P)\cdot u
        =
        (3e,0)\cdot(2,1)
        =
        6e.
    \]

    \item La máxima tasa de cambio de $f$ en $P$ se alcanza en la dirección del gradiente, es decir, tomando $u=\nabla f(P)=(3e,0)$. La derivada direccional en esa dirección es:
    \[
        f'(P;u)
        =
        \nabla f(P)\cdot u
        =
        (3e,0)\cdot(3e,0)
        =
        9e^2.\qedhere
    \]
\end{enumerate}
\end{respuesta}


%%%%%%%%%%%%%%%%%%%%%%%%%%%%%%%%%%%%%%%%%%
\item
Aplique la regla de la cadena para calcular las derivadas indicadas.
\begin{enumerate}[leftmargin=*]
    \item Calcule $\dfrac{dz}{dt}$ si
    $\quad
        z=x^2\sin(y),
        \quad
        x=e^t,
        \quad
        y=t^2.
    $\puntaje{3}

    \item Calcule $\dfrac{dz}{dt}$ si
    $
        \quad
        z=e^{xy},
        \quad
        x=u^2+v,
        \quad
        y=u-v^2,
        \quad
        u=\cos(t),
        \quad
        v=\sin(t).
    $\puntaje{5}
\end{enumerate}

\begin{respuesta}
\begin{enumerate}
    \item Aplicamos la regla de la cadena:
    \[
        \frac{dz}{dt}
        =
        \frac{\partial z}{\partial x}\cdot\frac{dx}{dt}
        +
        \frac{\partial z}{\partial y}\cdot\frac{dy}{dt}.
    \]
    Calculamos cada factor:
    \[
        \frac{\partial z}{\partial x}=2x\sen(y),\qquad
        \frac{\partial z}{\partial y}=x^2\cos(y),\qquad
        \frac{dx}{dt}=e^t,\qquad
        \frac{dy}{dt}=2t.
    \]
    Por tanto:
    \[
        \frac{dz}{dt}
        =
        2x\sen(y)\cdot e^t+x^2\cos(y)\cdot 2t.
    \]

    \item Aplicamos la regla de la cadena considerando la dependencia completa:
    \[
        \frac{dz}{dt}
        =
        \frac{\partial z}{\partial x}\cdot\frac{dx}{dt}
        +
        \frac{\partial z}{\partial y}\cdot\frac{dy}{dt},
    \]
    donde a su vez:
    \[
        \frac{dx}{dt}
        =
        \frac{\partial x}{\partial u}\cdot\frac{du}{dt}
        +
        \frac{\partial x}{\partial v}\cdot\frac{dv}{dt},
        \qquad
        \frac{dy}{dt}
        =
        \frac{\partial y}{\partial u}\cdot\frac{du}{dt}
        +
        \frac{\partial y}{\partial v}\cdot\frac{dv}{dt}.
    \]
    Calculamos:
    \[
        \frac{\partial z}{\partial x}=ye^{xy},\quad
        \frac{\partial z}{\partial y}=xe^{xy},\quad
        \frac{\partial x}{\partial u}=2u,\quad
        \frac{\partial x}{\partial v}=1,\quad
        \frac{\partial y}{\partial u}=1,\quad
        \frac{\partial y}{\partial v}=-2v,
    \]
    \[
        \frac{du}{dt}=-\sen(t),\qquad\frac{dv}{dt}=\cos(t).
    \]
    Entonces:
    \[
        \frac{dx}{dt}
        =
        2u(-\sen t)+\cos t
        =
        -2u\,\sen t+\cos t,
    \]
    y
    \[
        \frac{dy}{dt}
        =
        -\sen t+(-2v)\cos t
        =
        -\sen t-2v\cos t.
    \]
    Por lo tanto:
    \[
        \frac{dz}{dt}
        =
        e^{xy}\bigl[y(-2u\,\sen t+\cos t)+x(-\sen t-2v\cos t)\bigr].\qedhere
    \]
\end{enumerate}
\end{respuesta}


%%%%%%%%%%%%%%%%%%%%%%%%%%%%%%%%%%%%%%%%%%
\item
Sea $f(x,y)=e^{2x}\sen(y)$.
\begin{enumerate}[leftmargin=*]
    \item Calcule la matriz Hessiana de $f$ y evalúe es $\left(0,\dfrac{\pi}{2}\right)$.\puntaje{4}
    \item Escriba el polinomio de Taylor de orden 2 de $f$ alrededor de $a=\!\left(0,\dfrac{\pi}{2}\right)$.\puntaje{5}
    \item Use el polinomio obtenido para estimar $f\!\left(0.1,\;\dfrac{\pi}{2}+0.1\right)$.\puntaje{3}
\end{enumerate}


\begin{respuesta}
\begin{enumerate}
    \item Calculemos las derivadas parciales de segundo orden. Primero:
    \[
        \frac{\partial f}{\partial x}(x,y)
        =
        2e^{2x}\sen(y),
        \qquad
        \frac{\partial f}{\partial y}(x,y)
        =
        e^{2x}\cos(y).
    \]
    Luego:
    \[
        \frac{\partial^2 f}{\partial x^2}(x,y)
        =
        4e^{2x}\sen(y),
        \qquad
        \frac{\partial^2 f}{\partial x\partial y}(x,y)
        =
        2e^{2x}\cos(y),
        \qquad
        \frac{\partial^2 f}{\partial y^2}(x,y)
        =
        -e^{2x}\sen(y).
    \]
    Evaluamos en $\left(0,\dfrac{\pi}{2}\right)$:
    \[
        \frac{\partial^2 f}{\partial x^2}\!\left(0,\tfrac{\pi}{2}\right)
        =
        4,
        \qquad
        \frac{\partial^2 f}{\partial x\partial y}\!\left(0,\tfrac{\pi}{2}\right)
        =
        0,
        \qquad
        \frac{\partial^2 f}{\partial y^2}\!\left(0,\tfrac{\pi}{2}\right)
        =
        -1.
    \]
    Por lo tanto:
    \[
        H_f\!\left(0,\tfrac{\pi}{2}\right)
        =
        \begin{pmatrix}
            4 & 0\\
            0 & -1
        \end{pmatrix}.
    \]

    \item Usemos la expresión:
    \[
        F(x,y)
        =
        f(a)
        +
        \nabla f(a)\cdot((x,y)-a)
        +
        \frac{1}{2!}[(x,y)-a]^T H_f(a)[(x,y)-a].
    \]
    En este caso $a=\!\left(0,\dfrac{\pi}{2}\right)$. Calculamos:
    \[
        f(a)=e^0\sen\!\left(\tfrac{\pi}{2}\right)=1,
        \qquad
        \nabla f(a)=\left(2e^0\sen\!\left(\tfrac{\pi}{2}\right),\;e^0\cos\!\left(\tfrac{\pi}{2}\right)\right)=(2,0).
    \]
    Denotando $(x,y)-a=\!\left(x,\,y-\dfrac{\pi}{2}\right)$, entonces:
    \[
        F(x,y)
        =
        1
        +
        (2,0)\cdot\!\left(x,\,y-\tfrac{\pi}{2}\right)
        +
        \frac{1}{2}
        \begin{pmatrix}
            x & y-\tfrac{\pi}{2}
        \end{pmatrix}
        \begin{pmatrix}
            4 & 0\\
            0 & -1
        \end{pmatrix}
        \begin{pmatrix}
            x\\[4pt]
            y-\tfrac{\pi}{2}
        \end{pmatrix}.
    \]
    Por lo tanto, la aproximación cuadrática es:
    \[
        F(x,y)
        =
        1+2x+2x^2-\frac{1}{2}\!\left(y-\frac{\pi}{2}\right)^{\!2}.
    \]

    \item Usamos el polinomio obtenido para estimar $f\!\left(0.1,\;\dfrac{\pi}{2}+0.1\right)$:
    \[
        F\!\left(0.1,\,\tfrac{\pi}{2}+0.1\right)
        =
        1+2(0.1)+2(0.1)^2-\tfrac{1}{2}(0.1)^2
        =
        1+0.2+0.02-0.005
        =
        1.215.\qedhere
    \]
\end{enumerate}
\end{respuesta}


%%%%%%%%%%%%%%%%%%%%%%%%%%%%%%%%%%%%%%%%%%
\item
Dado el campo vectorial $F(x,y)=(x^2y+y^2x,\ x^2y^2)$:
\begin{enumerate}[leftmargin=*]
    \item Calcule $\mathrm{div}(F)$ y $\mathrm{rot}(F)$.\puntaje{4}
    \item Determine dos punto del plano donde $F$ tiene fuentes y dos donde tiene sumideros.\puntaje{4}
    \item Determine un punto donde $\mathrm{rot}(F)=0$.\puntaje{2}
\end{enumerate}


\begin{respuesta}
\begin{enumerate}
    \item Sea $F(x,y)=(f_1(x,y),f_2(x,y))$, donde:
    \[
        f_1(x,y)=x^2y+y^2x,
        \qquad
        f_2(x,y)=x^2y^2.
    \]
    La divergencia de $F$ es:
    \[
        \mathrm{div}(F)(x,y)
        =
        \frac{\partial f_1}{\partial x}(x,y)
        +
        \frac{\partial f_2}{\partial y}(x,y)
        =
        (2xy+y^2)+2x^2y.
    \]
    El rotacional de $F$ en el plano es:
    \[
        \mathrm{rot}(F)(x,y)
        =
        \frac{\partial f_2}{\partial x}(x,y)
        -
        \frac{\partial f_1}{\partial y}(x,y)
        =
        2xy^2-(x^2+2xy).
    \]

    \item Hay fuentes donde $\mathrm{div}(F)(x,y)>0$. Evaluamos en algunos puntos:
    \[
        \mathrm{div}(F)(1,1)=2(1)(1)+(1)^2+2(1)^2(1)=5>0,
    \]
    \[
        \mathrm{div}(F)(2,1)=2(2)(1)+(1)^2+2(2)^2(1)=13>0.
    \]
    Por tanto, $(1,1)$ y $(2,1)$ son fuentes. Hay sumideros donde $\mathrm{div}(F)(x,y)<0$:
    \[
        \mathrm{div}(F)(1,-2)=2(1)(-2)+(-2)^2+2(1)^2(-2)=-4+4-4=-4<0,
    \]
    \[
        \mathrm{div}(F)(2,-1)=2(2)(-1)+(-1)^2+2(2)^2(-1)=-4+1-8=-11<0.
    \]
    Por tanto, $(1,-2)$ y $(2,-1)$ son sumideros.

    \item Para determinar dónde $\mathrm{rot}(F)=0$, notamos que:
    \[
        \mathrm{rot}(F)(x,y)
        =
        2xy^2-x^2-2xy
        =
        x(2y^2-x-2y).
    \]
    En $x=0$ el rotacional se anula. Por ejemplo, en el punto $(0,1)$:
    \[
        \mathrm{rot}(F)(0,1)
        =
        0\cdot(2-0-2)
        =
        0.\qedhere
    \]
\end{enumerate}
\end{respuesta}


\end{preguntas}

\end{document}
